% GENERATED FILE. Edit canonical lessons, then run npm run generate:books.
% canonical-source-hash: fnv1a64:b99b85a000d40886
% canonical-lessons: ML-C470-petrol, ML-C470-indhanam, ML-C470-tayar, ML-C470-brekku, ML-C470-pancarakuka

\chapter{Fuel and Tyres}
\label{ch:ml-fuel-and-tyres}
\hlchaptermodality{470}

\begin{chapteropening}
\textbf{By the end of this chapter:} \emph{I can say petrol, fuel, a tyre, a brake and to get a puncture.}

Complete the last lesson of chapter 470: I can say petrol, fuel, a tyre, a brake and to get a puncture.
\end{chapteropening}

\section[\texorpdfstring{\ml{പെട്രോൾ} (peṭrōḷ)}{പെട്രോൾ (peṭrōḷ)}]{\ml{പെട്രോൾ} (peṭrōḷ) — petrol}

\label{lesson:ML-C470-petrol}



\begin{quote}
Before the new one: say the Malayalam for to make way, then the Malayalam for to cross over.
\end{quote}

\subsection*{You'll want to know: \ml{പെട്രോൾ}}
\textbf{\ml{പെട്രോൾ}} — \emph{peṭrōḷ} — \textquotedblleft{}petrol\textquotedblright{}.

\begin{quote}
\ml{പെട്രോളിന്റെ} \ml{വില} \ml{വീണ്ടും} \ml{കൂടി}.
\end{quote}

\begin{quote}
\emph{peṭrōḷinṟe vila vīṇṭuṁ kūṭi.}
\end{quote}

\begin{quote}
The price of petrol went up again.
\end{quote}

A petrol station is a \ml{പെട്രോൾ} \ml{പമ്പ്} in everyday Malayalam.

\begin{cousinweb}
\textbf{\ml{പെട്രോൾ}} (\emph{peṭrōḷ}) was borrowed from English \emph{petrol}, the fuel of cars and motorbikes.
\end{cousinweb}

\begin{grammarlens}[title={What you've built}]
One more word for getting around.
\end{grammarlens}

\subsection*{Your turn}
\begin{itemize}
\raggedright
  \item \emph{Say it:} \emph{peṭrōḷ}
  \item \emph{Say it:} \emph{peṭrōḷ}, once more
  \item \emph{Say it:} \emph{peṭrōḷinṟe vila vīṇṭuṁ kūṭi.} — The price of petrol went up again.
  \item \emph{Recall:} say the Malayalam for to make way, then the Malayalam for to cross over, then say \emph{peṭrōḷ} again
\end{itemize}

\begin{tcolorbox}[breakable,colback=teal!4,colframe=teal!35!black,title={Before you move on}]
What does \ml{പെട്രോൾ} mean? (\textquotedblleft{}Petrol\textquotedblright{}.) Say it once more.
\end{tcolorbox}
\section[\texorpdfstring{\ml{ഇന്ധനം} (indhanaṁ)}{ഇന്ധനം (indhanaṁ)}]{\ml{ഇന്ധനം} (indhanaṁ) — fuel}

\label{lesson:ML-C470-indhanam}



\begin{quote}
Before the new one: say the Malayalam for to cross over, then the Malayalam for petrol.
\end{quote}

\subsection*{You'll want to know: \ml{ഇന്ധനം}}
\textbf{\ml{ഇന്ധനം}} — \emph{indhanaṁ} — \textquotedblleft{}fuel\textquotedblright{}.

\begin{quote}
\ml{വിമാനത്തിന്} \ml{ധാരാളം} \ml{ഇന്ധനം} \ml{വേണം}.
\end{quote}

\begin{quote}
\emph{vimānattinŭ dhārāḷaṁ indhanaṁ vēṇaṁ.}
\end{quote}

\begin{quote}
A plane needs a lot of fuel.
\end{quote}

The general word: petrol, diesel, gas and firewood are all \ml{ഇന്ധനം}.

\begin{cousinweb}
\textbf{\ml{ഇന്ധനം}} (\emph{indhanaṁ}) is borrowed from Sanskrit \emph{indhana}, \textquotedblleft{}kindling, lighting\textquotedblright{}, and so what is burnt (Monier-Williams).
\end{cousinweb}

\begin{grammarlens}[title={What you've built}]
One more word for getting around.
\end{grammarlens}

\subsection*{Your turn}
\begin{itemize}
\raggedright
  \item \emph{Say it:} \emph{indhanaṁ}
  \item \emph{Say it:} \emph{indhanaṁ}, once more
  \item \emph{Say it:} \emph{vimānattinŭ dhārāḷaṁ indhanaṁ vēṇaṁ.} — A plane needs a lot of fuel.
  \item \emph{Recall:} say the Malayalam for to cross over, then the Malayalam for petrol, then say \emph{indhanaṁ} again
\end{itemize}

\begin{tcolorbox}[breakable,colback=teal!4,colframe=teal!35!black,title={Before you move on}]
What does \ml{ഇന്ധനം} mean? (\textquotedblleft{}Fuel\textquotedblright{}.) Say it once more.
\end{tcolorbox}
\section[\texorpdfstring{\ml{ടയർ} (ṭayar)}{ടയർ (ṭayar)}]{\ml{ടയർ} (ṭayar) — a tyre}

\label{lesson:ML-C470-tayar}



\begin{quote}
Before the new one: say the Malayalam for petrol, then the Malayalam for fuel.
\end{quote}

\subsection*{You'll want to know: \ml{ടയർ}}
\textbf{\ml{ടയർ}} — \emph{ṭayar} — \textquotedblleft{}a tyre\textquotedblright{}.

\begin{quote}
\ml{കാറിന്റെ} \ml{ഒരു} \ml{ടയർ} \ml{പഴയതാണ്}.
\end{quote}

\begin{quote}
\emph{kāṟinṟe oru ṭayar paḻayatāṇŭ.}
\end{quote}

\begin{quote}
One of the car's tyres is old.
\end{quote}

Plural: \ml{ടയറുകൾ}, \emph{ṭayaṟukaḷ}.

Plural: \ml{ടയറുകൾ}.

\begin{cousinweb}
\textbf{\ml{ടയർ}} (\emph{ṭayar}) was borrowed from English \emph{tyre}, the rubber ring around a wheel.
\end{cousinweb}

\begin{grammarlens}[title={What you've built}]
One more word for getting around.
\end{grammarlens}

\subsection*{Your turn}
\begin{itemize}
\raggedright
  \item \emph{Say it:} \emph{ṭayar}
  \item \emph{Say it:} \emph{ṭayar}, once more
  \item \emph{Say it:} \emph{kāṟinṟe oru ṭayar paḻayatāṇŭ.} — One of the car's tyres is old.
  \item \emph{Recall:} say the Malayalam for petrol, then the Malayalam for fuel, then say \emph{ṭayar} again
\end{itemize}

\begin{tcolorbox}[breakable,colback=teal!4,colframe=teal!35!black,title={Before you move on}]
What does \ml{ടയർ} mean? (\textquotedblleft{}A tyre\textquotedblright{}.) Say it once more.
\end{tcolorbox}
\section[\texorpdfstring{\ml{ബ്രേക്ക്} (brēkkŭ)}{ബ്രേക്ക് (brēkkŭ)}]{\ml{ബ്രേക്ക്} (brēkkŭ) — a brake}

\label{lesson:ML-C470-brekku}



\begin{quote}
Before the new one: say the Malayalam for fuel, then the Malayalam for a tyre.
\end{quote}

\subsection*{You'll want to know: \ml{ബ്രേക്ക്}}
\textbf{\ml{ബ്രേക്ക്}} — \emph{brēkkŭ} — \textquotedblleft{}a brake\textquotedblright{}.

\begin{quote}
\ml{അവൻ} \ml{പെട്ടെന്ന്} \ml{ബ്രേക്ക്} \ml{ചവിട്ടി}.
\end{quote}

\begin{quote}
\emph{avan peṭṭennŭ brēkkŭ caviṭṭi.}
\end{quote}

\begin{quote}
He suddenly stepped on the brake.
\end{quote}

You step on a brake in Malayalam: \ml{ബ്രേക്ക്} \ml{ചവിട്ടുക}.

\begin{cousinweb}
\textbf{\ml{ബ്രേക്ക്}} (\emph{brēkkŭ}) was borrowed from English \emph{brake}, the part that slows a vehicle.
\end{cousinweb}

\begin{grammarlens}[title={What you've built}]
One more word for getting around.
\end{grammarlens}

\subsection*{Your turn}
\begin{itemize}
\raggedright
  \item \emph{Say it:} \emph{brēkkŭ}
  \item \emph{Say it:} \emph{brēkkŭ}, once more
  \item \emph{Say it:} \emph{avan peṭṭennŭ brēkkŭ caviṭṭi.} — He suddenly stepped on the brake.
  \item \emph{Recall:} say the Malayalam for fuel, then the Malayalam for a tyre, then say \emph{brēkkŭ} again
\end{itemize}

\begin{tcolorbox}[breakable,colback=teal!4,colframe=teal!35!black,title={Before you move on}]
What does \ml{ബ്രേക്ക്} mean? (\textquotedblleft{}A brake\textquotedblright{}.) Say it once more.
\end{tcolorbox}
\section[\texorpdfstring{\ml{പഞ്ചറാകുക} (pañcaṟākuka)}{പഞ്ചറാകുക (pañcaṟākuka)}]{\ml{പഞ്ചറാകുക} (pañcaṟākuka) — to get a puncture}

\label{lesson:ML-C470-pancarakuka}



\begin{quote}
Before the new one: say the Malayalam for a tyre, then the Malayalam for a brake.
\end{quote}

\subsection*{You'll want to know: \ml{പഞ്ചറാകുക}}
\textbf{\ml{പഞ്ചറാകുക}} — \emph{pañcaṟākuka} — \textquotedblleft{}to get a puncture\textquotedblright{}.

\begin{quote}
\ml{വഴിയിൽ} \ml{വച്ച്} \ml{സൈക്കിളിന്റെ} \ml{ടയർ} \ml{പഞ്ചറായി}.
\end{quote}

\begin{quote}
\emph{vaḻiyil vaccŭ saikkiḷinṟe ṭayar pañcaṟāyi.}
\end{quote}

\begin{quote}
On the way the bicycle's tyre got a puncture.
\end{quote}

Past: \ml{പഞ്ചറായി}, \emph{pañcaṟāyi}.

English puncture plus \ml{ആകുക}, to become. The noun alone is \ml{പഞ്ചർ}.

\begin{cousinweb}
\textbf{\ml{പഞ്ചറാകുക}} (\emph{pañcaṟākuka}) is a Malayalam verb built from \ml{പഞ്ചർ}, which was borrowed from English \emph{puncture}, and \ml{ആകുക}, to become.
\end{cousinweb}

\begin{grammarlens}[title={What you've built}]
One more word for getting around.
\end{grammarlens}

\subsection*{Your turn}
\begin{itemize}
\raggedright
  \item \emph{Say it:} \emph{pañcaṟākuka}
  \item \emph{Say it:} \emph{pañcaṟākuka}, once more
  \item \emph{Say it:} \emph{vaḻiyil vaccŭ saikkiḷinṟe ṭayar pañcaṟāyi.} — On the way the bicycle's tyre got a puncture.
  \item \emph{Recall:} say the Malayalam for a tyre, then the Malayalam for a brake, then say \emph{pañcaṟākuka} again
\end{itemize}

\begin{tcolorbox}[breakable,colback=teal!4,colframe=teal!35!black,title={Before you move on}]
What does \ml{പഞ്ചറാകുക} mean? (\textquotedblleft{}To get a puncture\textquotedblright{}.) Say it once more.
\end{tcolorbox}
